\subsection{Banzhaf Power Index}

\content{
        \begin{definition} To calculate the Banzhaf Power Index: 
        \begin{itemize}
        \item List all winning coalitions.
        \item In each winning coalition, determine all players who are critical.
        \item Count up how many times each player is critical. 
        \item Convert these to fractions by dividing by the total number of times any
        player is critical. 
        \end{itemize}
        \end{definition}

        \begin{example} Calculate the Banzhaf Power Index for the three players
        in the following weighted voting system 
        $$ [10: 7, 6, 2]. $$
        \end{example}
}{% presentation
\vspace{3in}
}{% nopresentation
\vspace{3in}
}{% comments
}


\content{
        \begin{example} Calculate the Banzhaf Power Index for the following
        weighted voting system
        $$ [ 8: 6,3,2 ]. $$
        \end{example}
}{% presentation
\vspace{3in}
}{% nopresentation
\vspace{3in}
}{% comments
$P_1 = 3/5$, $P_2 = 1/5$ and $P_3 =1/5$.
}

\content{
        \begin{example} An executive board consists of a president (P) and three 
        vice-presidents (V1,V2,V3).  For a motion to pass it must have three yes 
        votes, one of which must be the president's. Find a weighted voting system 
        to represent this situation, then find the Banzhaf power index of the
        system. The list of all coalitions is given below
        \end{example}
}{% presentation
\begin{align*}
\{P_1\} \quad \{P_2\} \quad \{P_3\} \quad \{P_4\} \\
\{P_1,P_2\} \quad \{P_1,P_3\}\quad \{P_1,P_4\} \quad \{P_2,P_3\}\\
\{P_2, P_4\} \quad \{P_3, P_4\} \quad \{P_1, P_2, P_3\} \quad \{P_1, P_3,
P_4\}\\
\{P_1, P_2, P_4\}\quad \{P_2,P_3, P_4\}\quad \{P_1, P_2, P_3,P_4\}
\end{align*}
}{% nopresentation
\begin{align*}
\{P_1\} \quad \{P_2\} \quad \{P_3\} \quad \{P_4\} \\
\{P_1,P_2\} \quad \{P_1,P_3\}\quad \{P_1,P_4\} \quad \{P_2,P_3\}\\
\{P_2, P_4\} \quad \{P_3, P_4\} \quad \{P_1, P_2, P_3\} \quad \{P_1, P_3,
P_4\}\\
\{P_1, P_2, P_4\}\quad \{P_2,P_3, P_4\}\quad \{P_1, P_2, P_3,P_4\}
\end{align*}
\vspace{1in}
}{% comments
One example of this is $[8: 4, 2, 2, 2]$. The Banzhaf power index of this
weighted voting system is
$$ P_1: 2/5\: P_2: 1/5\: P_3: 1/5\: P_4: 1/5 $$
}


\content{
        \begin{example} Compute the Banzhaf power index for the following
        weighted voting system. 
        $$ [16: 7,6,3,3,2 ]. $$
        The winning coalitions are listed below: 
        \begin{align*} 
        \{P_1, P_2, P_3\}, \{P_1, P_2, P_4\}, \{P_1, P_2, P_3, P_4\}, \\
        \{P_1, P_2, P_3, P_5\}, \{P_1, P_2, P_4, P_5\}, \\
        \{P_1, P_2, P_3, P_4, P_5\}
        \end{align*}
        \end{example}
}{% presentation
\vspace{1in}
}{% nopresentation
\vspace{1in}
}{% comments
Power index for this example is 
$$ P_1: 3/8\quad P_2: 3/8\quad P_3: 1/8 $$
$$P_4: 1/8\quad P_5: 0 $$
}

